\documentclass[12pt]{book}
\usepackage{amsmath}
\usepackage{amssymb}
\usepackage{geometry}
\geometry{margin=1in}

\title{Volume 08: Cosmology and Large-Scale Structure \\ Book 01: Cosmological Structure from Flux Topology}
\author{James C. Harvey}
\date{December 2025}

\begin{document}
\maketitle
\tableofcontents

\chapter{Cosmological Structure as Flux Geometry}

\section*{Abstract}
Large-scale structure is the macroscopic topology of flux in the matrix. Galactic rotation, disk
occlusion, and cosmic filaments are coherent regimes of the same coupling equation.

\section{Disk Occlusion and Rotation}
Disk occlusion modifies the pressure field:
\begin{equation}
\Pi(\mathbf{r}) = \int_{\Omega} I_{\text{CMB}}(\hat{\mathbf{n}})\left[1 - E(\mathbf{r},\hat{\mathbf{n}})\right]\, d\Omega.
\end{equation}
Rotation curves follow from gradients in $\Pi$ across disk geometry.

\section{Topology of Large-Scale Flux}
Flux pathways create filamentary structure by minimizing gradient cost across the matrix.

\chapter{CMB as Cosmological Boundary Condition}

\section*{Abstract}
The CMB is the universal influx and sets the cosmological pressure baseline. It is the boundary condition
for large-scale structure formation.

\section{Baseline Pressure}
\begin{equation}
P_{\text{CMB}} = 2.036 \times 10^{-2}\ \text{Pa}.
\end{equation}

\chapter{Expansion as Flux Redistribution}

\section*{Abstract}
Cosmic expansion is a redistribution of flux pathways rather than creation of space. It is the global
reconfiguration of pressure gradients.

\section{Scale Relation}
\begin{equation}
\frac{d}{dt}(\nabla \Pi) \neq 0,
\end{equation}
indicating evolving large-scale pressure topology.

\chapter{Galactic Rotation and Occlusion}

\section*{Abstract}
Rotation curves are the result of disk occlusion patterns. The effective field is the integral of
occlusion across the disk geometry.

\section{Rotation Condition}
\begin{equation}
v^2(r) = r\, \frac{d\Phi_G}{dr},
\end{equation}
with $\Phi_G$ derived from occlusion-weighted pressure.

\chapter{Filaments and Voids}

\section*{Abstract}
Filaments are flux highways; voids are low-gradient cavities. Their distribution follows the minimization
of global pressure gradients.

\section{Topology Rule}
\begin{equation}
\delta \Pi \rightarrow \text{min} \quad \Rightarrow \quad \text{filament network}.
\end{equation}
\end{document}
